\documentclass[12pt,a4paper]{article}
\usepackage[utf8]{inputenc}
\usepackage[T1]{fontenc}
\usepackage[spanish]{babel}
\usepackage[margin=3cm]{geometry}
\usepackage{setspace}
\onehalfspacing

\title{La libertad en Sartre: condena y proyecto}
\author{Julián Martínez Ortega\\\small Introducción a la Filosofía Contemporánea}
\date{Octubre de 2026}

\begin{document}
\maketitle

\section{Introducción}
Pocas frases han tenido tanta fortuna en la filosofía del siglo \textsc{xx} como la
afirmación de Jean-Paul Sartre según la cual ``el hombre está condenado a ser
libre''.\footnote{J.-P. Sartre, \emph{El existencialismo es un humanismo}, 1946.}
En este ensayo examino el sentido de esa condena y la relación que guarda con la noción
de proyecto.

\section{La existencia precede a la esencia}
Para Sartre, a diferencia de un cortapapeles, el ser humano no tiene una naturaleza
fijada de antemano:
\begin{quote}
El hombre empieza por existir, se encuentra, surge en el mundo, y después se define. El
hombre, tal como lo concibe el existencialista, si no es definible, es porque empieza
por no ser nada.
\end{quote}
De aquí se sigue que no hay excusas: no podemos apelar a una naturaleza humana, a Dios o
a las circunstancias para justificar nuestras elecciones.

\section{Angustia, abandono y desesperación}
Sartre describe tres actitudes que acompañan a la libertad:
\begin{enumerate}
  \item la \emph{angustia}, porque al elegir para mí elijo una imagen del hombre;
  \item el \emph{abandono}, porque no hay valores dados de antemano;
  \item la \emph{desesperación}, porque solo cuento con lo que depende de mi voluntad.
\end{enumerate}

\section{Objeciones}
Merleau-Ponty objetó que Sartre ignora el peso de la situación: la libertad siempre está
encarnada y sedimentada en hábitos.\footnote{M. Merleau-Ponty, \emph{Fenomenología de la
percepción}, 1945, cap.~III.} Esta crítica obliga a matizar la tesis sin abandonarla.

\section{Conclusión}
La condena a la libertad no es un castigo sino una descripción de nuestra condición: somos
el proyecto que elegimos ser, y somos responsables de él.

\end{document}
